% Синтетический кейс: то, что на круге tex -> docx -> tex -> docx менялось от круга к кругу
% (стресс-прогон 2026-11-09): пустые разделители \left. \right., норма \lVert, align со строкой
% \notag и строкой с cases, теорема внутри пункта списка, формула с меткой в сноске.
\documentclass[12pt]{article}
\usepackage[T2A]{fontenc}
\usepackage[utf8]{inputenc}
\usepackage[russian]{babel}
\usepackage{amsmath,amsthm}

\newtheorem{lemma}{Лемма}

\begin{document}

\section{Разделители}

Норма $\left\lVert v \right\rVert$, производная $\left.\frac{df}{dx}\right|_{x=0}$
и скобка $\left\{ a \right.$.

\begin{align}
  f(x) &= x^2 + 2x + 1 \label{eq:f}\\
       &= (x+1)^2 \notag\\
  g(x) &= \begin{cases} 0, & x < 0,\\ x, & x \ge 0 \end{cases} \label{eq:g}\\
  h(x) &= \left\lVert x \right\rVert. \label{eq:h}
\end{align}

\section{Списки и сноски}

\begin{itemize}
  \item Первый пункт:
  \begin{lemma}\label{lem:a}
  Лемма в пункте списка.
  \end{lemma}
  \item Второй пункт, см. лемму~\ref{lem:a}.
\end{itemize}

Текст со сноской\footnote{Сноска с формулой
\begin{equation}\label{eq:fn}
  a = b
\end{equation}
и текстом после.}. Ссылки \eqref{eq:fn}, \eqref{eq:f}--\eqref{eq:h} и
\begin{equation}\label{eq:last}
  c = d.
\end{equation}

\end{document}
